\begin{pastexam}[label={px:gw3-110-ds-01}]{關務三等資訊處理 110 年 第一題}
（申論）

大學生只剩 5 天準備 4 科 X$_1$, X$_2$, X$_3$, X$_4$，估計的成績點數如下表所示，每 1 科準備至少 1 天，使用窮舉法
（exhaustive search）有幾種可能？如何最適化？最適成績 s =？（20 分）
\begin{center}
\begin{tabular}{|l|c|c|c|c|}
\hline
天數 \hfill 科目 & X$_1$ & X$_2$ & X$_3$ & X$_4$ \\
\hline
1（天） & 3 & 5 & 2 & 4 \\
\hline
2 & 5 & 6 & 4 & 4 \\
\hline
3 & 6 & 8 & 7 & 5 \\
\hline
\end{tabular}
\end{center}
\end{pastexam}

\begin{pastsoln}
答案：$4$ 種可能；用動態規劃逐科決定天數；最適成績 $s = 16$（X$_1$ 準備 $2$ 天、其餘各 $1$ 天，或 X$_3$ 準備 $2$ 天、
其餘各 $1$ 天）。

\textbf{窮舉。} 設 X$_i$ 準備 $d_i$ 天，$d_i \ge 1$，$d_1 + d_2 + d_3 + d_4 = 5$。 四科至少各 $1$ 天已用掉 $4$ 天，只剩
$1$ 天可以加給某一科，所以只有 $4$ 種分法（一般地，$5$ 拆成 $4$ 個正整數的有序拆法有 $\binom{5-1}{4-1} = \binom{4}{3} = 4$ 種）：
\begin{center}
\begin{tabular}{cl}
\toprule
$(d_1, d_2, d_3, d_4)$ & 成績 \\
\midrule
$(2, 1, 1, 1)$ & $5 + 5 + 2 + 4 = 16$ \\
$(1, 2, 1, 1)$ & $3 + 6 + 2 + 4 = 15$ \\
$(1, 1, 2, 1)$ & $3 + 5 + 4 + 4 = 16$ \\
$(1, 1, 1, 2)$ & $3 + 5 + 2 + 4 = 14$ \\
\bottomrule
\end{tabular}
\end{center}
（若寫成程式，讓每科的天數各在 $1 \sim 3$ 之間變動，共試 $3^4 = 81$ 種組合，其中只有這 $4$ 種剛好 $5$ 天。 表中多準備
一天成績都不會下降，所以把 $5$ 天用完一定不吃虧；若允許不用完，只多出全部 $1$ 天的一種，成績 $14$，不影響答案。）

\textbf{最適化：動態規劃。} 窮舉在科目或天數多時會爆炸（$m$ 科、每科 $1 \sim t$ 天要試 $t^m$ 種），改用動態規劃
（資結書第~18.1~節，這是「資源分配」型的題目）。 令 $p_i(k)$ 是 X$_i$ 準備 $k$ 天的成績，$f(i, d)$ 是「前 $i$ 科共用
$d$ 天、每科至少 $1$ 天」的最高成績：
\[
f(1, d) = p_1(d),\qquad
f(i, d) = \max_{1 \le k \le 3,\ k \le d - (i-1)} \bigl\{ f(i-1, d-k) + p_i(k) \bigr\}.
\]
最優子結構：最佳分配裡，前 $i - 1$ 科用掉 $d - k$ 天的那部分，一定也是「$i - 1$ 科、$d - k$ 天」的最佳分配，否則換成
更好的就能提高總分。 答案是 $f(4, 5)$。 由於後面每科至少要留 $1$ 天，前 $i$ 科最多用 $5 - (4 - i)$ 天，只需填下表：
\begin{center}
\begin{tabular}{lcccc}
\toprule
 & $d = 1$ & $d = 2$ & $d = 3$ & $d = 4$ \\
\midrule
$f(1, d)$ & 3 & 5 & & \\
$f(2, d)$ & & $3 + 5 = 8$ & $\max(5 + 5,\ 3 + 6) = 10$ & \\
$f(3, d)$ & & & $8 + 2 = 10$ & $\max(10 + 2,\ 8 + 4) = 12$ \\
\bottomrule
\end{tabular}
\end{center}
\[
f(4, 5) = \max\bigl( f(3, 4) + p_4(1),\ f(3, 3) + p_4(2) \bigr) = \max(12 + 4,\ 10 + 4) = 16.
\]
回溯：$f(4, 5)$ 取自 $f(3, 4) + p_4(1)$，X$_4$ 準備 $1$ 天；$f(3, 4) = 12$ 有兩種來源——$f(2, 3) + p_3(1)$（X$_3$ $1$ 天，
而 $f(2, 3) = 10$ 來自 X$_1$ $2$ 天、X$_2$ $1$ 天）或 $f(2, 2) + p_3(2)$（X$_3$ $2$ 天，X$_1$、X$_2$ 各 $1$ 天）。 所以
最適解有兩個：$(2, 1, 1, 1)$ 和 $(1, 1, 2, 1)$，最適成績 $s = 16$，和窮舉的結果一致。

一般情形 $m$ 科、$D$ 天、每科最多 $t$ 天，表格有 $O(mD)$ 格、每格試 $O(t)$ 種，時間 $O(mDt)$，遠小於窮舉的 $t^m$。
\end{pastsoln}
